\documentclass{report}
\usepackage{amsmath,amssymb}
\setlength{\parindent}{0mm}
\setlength{\parskip}{1em}
\begin{document}
\begin{center}
\rule{6in}{1pt} \
{\large Daniela Calvetti \\
{\bf Statistically inspired preconditioned iterative solvers for ill-posed linear systems}}

Case Western Reserve University \\ MAMS Department \\ 10900 Euclid Ave \\ Cleveland \\ OH 44118
\\
{\tt dxc57@case.edu}\\
Erkki Somersalo\\
Jamie Prezioso\\
Francesca Pitolli\\
Barbara Vantaggi\end{center}

The solution of ill-posed linear systems via iterative methods is a
versatile tool to account for qualitative features that are believed to
characterize the solution when recasting the problem in a probabilistic
framework where every unknown is modeled as a random variable. In this
talk we show how to use a hierachical model and anatomical features of
the brain to design a class of right preconditioners for the solution of
the severely underdetermined inverse problem of magnetoenchephalography.
These preconditioners depend on a set of parameters that are estimated
together with the unknown of primary interest. In the general setup, the
parameters can be interpreted as the variances of the entries of the
unknown, with a higher parameter value indicating that the corresponding
solution entry may have a value significantly different from the nearby
entries.
The algorithm for estimating the unknown and the parameters in the
preconditioner is an iterative alternating scheme (IAS) which updates
alternatively one of the two set of unknowns, while keeping the other
fixed. The IAS algorithm can be computationally very efficient: the
updating of the unknown requires the solution of a least squares problem,
while an explicit formula can be used to update the parameters when they
are assumed to follow a Gamma distribution. Moreover, the IAS algorithm
is globally convergent and it can be shown that one of the two parameters
of the Gamma distribution can be tuned to promote sparsity. A
statistically motivated criterion for stopping the iteration in the
solution of the least squares problems is also presented.


\end{document}
